\documentclass[12pt]{report}
\usepackage{graphicx}
\graphicspath{ {images/} }

\title{
{Implementation of publihser / subscriber inter process 
communication model at kernel level}
\newline

{\large Harcourt Butler Technical University, Kanpur}
\newline

{\includegraphics{hbtu_logo.png}}
}
\author{
	Ayush Kumar Shakya
	\and
	Arunima Verma
	\and
	Ashutosh Singh
	\and
	Hemant Singh
}
\date{\today}

\begin{document}
	\maketitle 
	
	\chapter*{Certificate}
	
	This is to certify that 
    \textbf{Ayush Kumar Shakya} (Roll no. 190108017),
	\textbf{Arunima Verma} (Roll no. 190108014),
	\textbf{Ashutosh Singh} (Roll no. 190108015),
	\textbf{Hemant Singh} (Roll no. 190108024),
	students of INFORMATION TECHNOLOGY, Harcourt Butler Technical University, 
	Kanpur are working on their project under my guidance.
	They have a desire for learning and I wish success in their future project work.
	
	\vfill
	\textit{Project guided by}
	\newline
	\textbf{Dr. Prabhat Verma}
	
	\addcontentsline{toc}{chapter}{Certificate}

	\chapter*{Abstract}
	
	Inter process communication (IPC) protocols must be provided by multitasking systems 
	in order to allow communication between processes with varying levels of privilege. 
	
	These mechanisms may be implemented in user space or at the application level, 
	however doing so at the application level has performance concerns because 
	it must make numerous syscalls, which add a lot of overhead. IPC at the kernel level 
	can decrease overheads and improve security. 
	
	Additionally, current kernel-based IPC models like Message Queues 
	require coordination between processes in order to communicate, which might have 
	overheads \cite{citation01} and prevent anonymous communication. 
	
	We believe that each of those problems can be resolved by a kernel level 
	publisher subscriber model.

	
	\addcontentsline{toc}{chapter}{Abstract}
	\tableofcontents
	

	\chapter{Introduction}
	In a modern system we have hundreds of processes running at any time and they mostly 
	communicate using application level IPC mechanisms.
	
	The majority of the time, processes choose not to use any communication techniques 
	and often build everything into a single programme, inflating the programmes 
	into large monoliths that are challenging to debug.
	
	Software can be made more reliable and scalable, perhaps across several devices, 
	if it is created as small microservices that utilise effective IPC mechanisms.
	
	\section{Why publisher-subscriber pattern ?}
    Publish/Subscribe messaging is an asynchronous service-to-service communi-
    cation method used in serverless and microservices architectures. Publishers
    and subscribers are decoupled from one another so that publishers do not
    know the destination of the information that they send, and subscribers do
    not know the source of the information that they receive.

    It has following benefits
	\begin{enumerate}
		\item Decoupled/Loosely systems
		\item Real\-time communication
		\item Ease of development
	\end{enumerate}
	
	\section{IPC Model}
	\begin{enumerate}
		\item Publishing
		\item Subscription
		\item Notification
	\end{enumerate}
	
	\section{Publishing}
	Any process may send any number of messages on a particular \textbf{topic} or 
	multiple topics, topics may be represented by an integer.

	some of the topics may be reserved for system use so programs listening to those
	topics can be sure that no malicious program can messages on reserved topics.
	
	This may be implemented as a syscall at kernel level.
				 	
	\section{Subscription}
	Processes can subscribe to one or more than one topics and register a callback
	function to be called when they receive a notification.
	kernel may store these processes in a hash table where \textit{key} is topic 
	and \textit{value} is a linked list or process identifiers or PIDs.
	This registration maybe implemented as a syscall.c
		
	\section{Notification}
	When a notification arrives on a particular topic, it is added to a queue,
	and kernel dequeues topics one at a time and sends a copy of message to each 
	process \cite{citation04}.
	
	We opted for copying for messages rather than shared memory because
	modern CPUs are quite good at copying small data and having a copy rather than
	shared resourse will help prevent any concurrency related issues such as deadlocks.	

    \section{Sharing information}
    \subsection{Shared memory}
    Inter process communication by shared memory is a concept where two or
    more process can access the common memory and communication is done
    through this shared memory where changes made by one process can be
    viewed by another process.

    For applications that exchange large amounts of data, shared memory is
    far superior to message passing techniques like message queues, which require
    system calls for every data exchange.

    To use shared memory, we have to perform two basic step
    \begin{enumerate}
		\item Request a memory segment that can be shared between processes to the operating system.		
        \item associate a part of that memory or the whole memory with the address space of the calling process.
	\end{enumerate}
    
    \subsection{Message Passing}
    Message passing can be used as a more process-oriented approach to synchronization 
    than the ”data-oriented” approaches used in providing mutual
    exclusion for shared resources.

    Message passing model allows multiple processes to read and write data
    to the message queue without being connected to each other. Messages are
    stored on the queue until their recipient retrieves them. Message queues are
    quite useful for interprocess communication and are used by most operating
    systems.

    We opted for message passing rather than shared memory because in
    shared memory we have to coordinate a lot between processes which has an
    overhead.

    Furthermore modern processors are very good at copying and moving
    small amount of data

	\section{Cleanup}
	When is process is no longer running it may be removed from the hash table
	so if the same process identifer is assigned to some different process it 
	doesn't receive messages it did not subscribe to.
    
	\chapter{Related Work}
	\section{Apache Kafka}
	Apacha Kafka is distributed event broker that provides publisher-subscriber 
	IPC mechanism at large distributed scale.It more or less achives the functionality 
	that we want to achive.
	
	It is built using Java and Scala and is JVM based therefore quite resourse intensive
	and therefore not suitable for performance critical workloads.
	
	\section{D-Bus}
	D-Bus or Desktop bus is a system bus that provides communication between processes
	using a single system bus it aims to provide standarization communication for 
	linux services.
	
	\subsection{Operating modes} 
	\begin{enumerate}
		\item Request response
		\item Publish subscribe
	\end{enumerate} 
	
	It implemented in user space and therefore cannot be used by kernel specially during
	booting when it needs to stacrt multiple services and a dbus like functionality would
	simplify and speed up the booting process.
	
	There proposals to add dbus like functionality at linux kernel level but nothing
	concrete is merged yet\cite{citation03}.	
	
    \chapter{Implementation}
		
	\section{Kernel}
	To reduce project scope and resource availability, 
	we have developed a 64 bit x86 kernel written in C and will concentrate 
	on optimising our ipc mechanisms for x86 platform only.
	
	The kernel we developed includes a bootloader that loads the kernel and initializes virtual memory
	as well switches to long mode (64 bit mode).
	
	Additionally, the kernel contains several utility functions for both the kernel and userspace programmes, 
	as well as common data structures for heap allocation. 
	
	In addition, the kernel configures the hardware timer, keyboard, and other configurations for CPU interrupts 
	from hardware and software.

    \section{Bootloader}
    \subsection{Loading Kernel}
    In the x86 architecture, the bootup procedure looks for a bootloader on
    the first sector of the booting device (a USB flash disk, hard drive, image
    file, etc.). This is automatically loaded and executed during the system
    booting process. The bootloader is responsible for loading the rest of the
    OS from the booting device into memory. The bootloader cannot rely upon
    any functionality from the OS, such as OS calls for disk I/O, or linking an
    object file against static libraries. However, the x86 architecture requires the
    presence of the Basic Input Output System (BIOS), which provides minimal
    terminal, keyboard and disk support. Your bootloader will use these BIOS
    calls to load the kernel from disk into memory, set up a stack space, and then
    switch control to the kernel. At this point, the OS begins running.

    \subsection{Memory Protection}
    Segmentation provides a mechanism of isolating individual code, data, and   
    stack modules so that multiple programs (or tasks) can run on the same processor 
    without interfering with one another.

    Whatever memory address you came across inside a process is actually a logical address. 
    This logical address then converts into linear address using segmentation, then into physical
    address(Address on RAM) using paging. Segmentation is a mechanism for
    dividing the linear address space into smaller protected address spaces called
    segments.

    Paging is a different way of accessing physical memory. Instead of memory 
    being divided into segments it is divided into pages and frames. 
    
    x86 processors normally use 3 or 4 level paging with most recent ones being able
    to use even 5 level paging systems. Using paging allows us to create arbitrary
    memory layouts and better utilize memory. It also allows us to implement
    more complex scenarios where each process has its own address space.

 
	\section{Process Managment and Scheduler}
	We plan to introduce a process managment and round robin scheduler
	to our kernel.
	
	In order to avoid linking-related activities, we have chosen to employ dummy programmes that 
	our kernel would interpret.
	
	\section{Program}
	A program for our kernel will be a static list of instructions with a fix length
	that will be interpeted instruction by instruction
	
	\begin{verbatim}
	typedef struct Program {
    	int data[MAX_PROGRAM_LEN];
    	int len;
    	int ip;
	} Program;
	\end{verbatim}
	
	data will contain instructions for interpretation as well as any input for those instructions
	\textit{ip} will contain an index that points to the next instruction.
	
	\subsection{Process management}
	The state of the process manager is kept behind a mutex so we don't get into a situation
	where our scheduler updates the state while we're executing our processes.
 
	\begin{verbatim}
	typedef struct
	{
		int current_pid;
		uint32_t mutex_lock;
	} ProcessState;
	static Program processes[256];
	static int pid_count = 0;
	static ProcessState state = {
   	    .current_pid = -1,
   	    .mutex_lock = MUTEX_UNLOCKED,
	};
	\end{verbatim}
	
	The scheduler will simply run in an infinite loop run the currently active task
	move the instruction pointer for the current task to the next instruction and wait for
	a certain amount of time to make sure all processes get equal amount of executuion time.
	
	A CPU timer interrupt will periodically give control back to the kernel from the
	scheduler and we move to the next available task.

    \section{Achieving multitasking}
    Now we develop two functions \textit{scheduler\_addtask} to add a new task/thread
    to schedule and manage another function \textit{scheduler\_start} that takes control from
    the kernel and starts executing the available tasks registered through \textit{scheduler\_addtask}.

    \begin{verbatim}
    void scheduler_addtask(Program *program)
    {
        int64_t pid = pid_count++;
        program->pid = pid;
        processes[pid] = program;
    }

    void scheduler_start()
    {
        while(1) 
        {
            spin_lock(&state.mutex_lock);
            if (state.current_pid > -1)
            {
                Program *p = processes[state.current_pid];
                program_step(p);
            }
            spin_unlock(&state.mutex_lock);
        }
    }
    \end{verbatim}
    
    \begin{verbatim}
        
    \end{verbatim}

    \chapter{Digging Deeper into Threads}
    Similar to JVM (Java virtual machine), we have implemented a virtual machine coupled with 
    our scheduler that loads the programs and interprets a low level byte code designed to be 
    platform independent and portable.

    Byte code  can be compiled from a higher level language such as C/C++.

    Initially byte code is loaded from disk into memory along with all its required static and 
    dynamic data and the current state of the thread is initialized i.e. setting the instruction 
    pointer to zero and assigning a process id.

    Then with the help of a hardware timer, the scheduler chooses a thread to run next and 
    mutex locks the current scheduler state till we execute one instruction in order to 
    prevent hardware timer to not change process during the execution of an instruction 
    i.e All instructions are atomic.

    Then each instruction is then passed to virtual machine that interprets the relevant 
    instructions and/or call system calls

    \section{Program Instructions}
    List of some instructions as expressed as C style enum
    \begin{verbatim}
    enum 
    {
    // Marks the start of your program
    PROGRAM_START = 0,
    
    // Does nothing
    PROGRAM_NOOP,
    
    // usage: [PROGRAM_LOAD_IMMEDIATE, REGISTER_A, 100]
    // it load the data into register A
    PROGRAM_LOAD_IMMEDIATE,

    // usage: [PROGRAM_ADD REGISTER_A, REGISTER_B] 
    // moves the data of register B into A
    PROGRAM_MOV,
  
    // usage: [PROGRAM_ADD REGISTER_A, REGISTER_B] 
    // it adds values in A and B together and stores them in A
    PROGRAM_ADD,

    // usage: [PROGRAM_INC REGISTER_A, amount] 
    // increments the value stored in register by amount
    PROGRAM_INC,

    // usage [PROGRAM_PRINT REGISTER_A]
    // it prints value A to screen
    PROGRAM_PRINT, 

    // usage [PROGRAM_SYSCALL, SYSCALL_ID, REGISTER_A]
    // calls syscall with the given id and stores the result in
    // REGISTER_RESULT
    PROGRAM_SYSCALL,

    // usage [PROGRAM_SYSCALL, SYSCALL_ID, SYSCALL_DATA]
    // calls syscall with the given id and stores the result in
    // REGISTER_RESULT
    PROGRAM_SYSCALL_IMMEDIATE,

    // usage [PROGRAM_TEST, REGISTER_A, REGISTER_B]
    // sets zero flag if REGISTER_A == REGISTER_B
    PROGRAM_TEST,

    // usage [PROGRAM_TEST, REGISTER_A, VALUE]
    // sets zero flag if REGISTER_A == VALUR
    PROGRAM_TEST_IMMEDIATE,

    // resets flags
    PROGRAM_RESET,

    // usage [PROGRAM_JUMP, ip]
    // jumps to instruction
    PROGRAM_JUMP,
  
    // usage [PROGRAM_JUMP_IF, ip]
    // jumps to instruction zero flag is set
    PROGRAM_JUMP_IF,

    // usage [PROGRAM_MOV_MEMORY_START, REGISTER_A]
    // load the memory base to register_a
    PROGRAM_LOAD_MEMORY_BASE,

    // usage [PROGRAM_DEREF, REGISTER_A, REGISTER_B]
    // deref the memory address stored in b and stores
    // the value in A
    PROGRAM_DEREF,

    // Marks the end of your program as well base base of 
    // of fixed memory base
    PROGRAM_END
    }
    \end{verbatim}

    \pagebreak
    
    \section{Hello World example}
    \begin{verbatim}
        PROGRAM_START,
        PROGRAM_LOAD_MEMORY_BASE, REGISTER_B,

        // 3
        PROGRAM_DEREF, REGISTER_C, REGISTER_B,
        PROGRAM_SYSCALL, SYSCALL_PUTCHAR, REGISTER_C,
        
        PROGRAM_TEST_IMMEDIATE, REGISTER_C, 0,
        PROGRAM_JUMP_IF, 19,
        PROGRAM_INC, REGISTER_B, 8,

        PROGRAM_JUMP, 3,

        // 19
        PROGRAM_NOOP,
        PROGRAM_END,
        'H','e','l','l','o',' ','W','o','r','l','d','\n',0,
    \end{verbatim}

    \chapter{Implementing Inter process communication}
	TODO: @arunima @hemant

 
	\pagebreak
	\begin{thebibliography}{9}
	\bibitem{citation01}
		S. L. Mirtaheri, E. M. Khaneghah and M. Sharifi, 
		A Case for Kernel Level Implementation of Inter Process Communication Mechanisms,
		2008 3rd International Conference on Information and Communication Technologies: From Theory to Applications, 
		2008, pp. 1-7, doi: 10.1109/ICTTA.2008.4530360.
	
	\bibitem{citation02}
		R. Rajkumar, M. Gagliardi and Lui Sha, 
		The real-time publisher/subscriber inter-process communication model for distributed real-time systems: design and implementation, 
		Proceedings Real-Time Technology and Applications Symposium, 
		1995, pp. 66-75, doi: 10.1109/RTTAS.1995.516203. 

	\bibitem{citation03}
		David Herrmann
		Bus IPC, https://lists.linuxfoundation.org/pipermail/ksummit-discuss/2016-July/003047.html
		2016
		
	\bibitem{citation04}
		Jochen Liedtke. 1993. 
		Improving IPC by kernel design. In Proceedings of the fourteenth ACM symposium on Operating systems principles (SOSP '93). 
		Association for Computing Machinery, 
		New York, NY, USA, 175–188. 
		https://doi.org/10.1145/168619.168633

    \bibitem{citation05}
        https://www.ijesi.org/papers/Vol(7)i5/Version-3/I0705035155.pdf

    \bibitem{citation07}
        https://www.researchgate.net/publication/232625032\_A\_Survey\_of
        \_Virtual\_Machine\_System\_Current\_Technology\_and\_Future\_Trends

    \bibitem{citation08}
    https://www.researchgate.net/publication/353958262\_A\_RESEARCH\_SURVEY\_ON
    \_VIRTUAL\_MACHINES\_FOR\_EFFICIENT\_COMPUTING\_IN\_OPERATING\_SYSTEMS

    \bibitem{citation09}
    https://iopscience.iop.org/article/10.1088/1742-6596/1950/1/012027/pdf
 
    \bibitem{citation10}
    https://compas.cs.stonybrook.edu/~nhonarmand/courses/fa14/cse506.2/papers/chen01.pdf

	\end{thebibliography}
	\addcontentsline{toc}{chapter}{Bibliography}

\end{document}
